\section{Ideas extras creativas}

\subsection{Velocidad de escape local durante el encuentro}

Calcular a lo largo del tiempo la velocidad de escape local del sistema Tierra--Apophis,
\[
v_{esc} = \sqrt{\frac{2 G M_{\oplus}}{r}},
\]
y compararla con la velocidad relativa de Apophis. Mostrar que en ningún momento $v_{rel} < v_{esc}$ y explicar por qué el asteroide no queda capturado.

\subsection{Comparar con la segunda ley de Kepler}

Calcular numéricamente el área barrida por el vector relativo Sol--Apophis en intervalos de tiempo iguales. Verificar que dichas áreas son iguales en la aproximación de 2 cuerpos y estudiar cómo esta ley se viola ligeramente cuando se incluyen perturbaciones planetarias.

\subsection{Energía de la hipérbola geocéntrica}

Calcular la energía específica geocéntrica,
\[
\mathcal{E}_{geo} = \frac{v^2}{2} - \frac{G M_{\oplus}}{r},
\]
a lo largo de la trayectoria de Apophis. Mostrar que permanece positiva, como corresponde a una trayectoria hiperbólica no capturada, y que alcanza un mínimo cerca del instante de máxima aproximación.

\subsection{Posición en el cielo desde Colombia}

Convertir las coordenadas eclípticas de Apophis en el momento del encuentro a ascensión recta y declinación. Determinar en qué constelación estará y evaluar si será visible desde Medellín durante la noche del 13 de abril de 2029.